\section{Routing}
\label{sec:routing}

\begin{lede}
How a robot works out a path to its next waypoint, and when that path is
recomputed.
\end{lede}

By default \code{Router.route} is \code{graph.shortest\_route}
(\Cref{sec:graph-distances}) --- plain BFS descent. When
\param{direction\_aware\_routing} is enabled \emph{and}
\param{direction\_control} is \code{"aisle"}, it instead runs a weighted
A$^\ast$ using the BFS distance map as an admissible and consistent heuristic:
\begin{align}
  \mathrm{penalty}(u, v) &=
  \begin{cases}
    0 & a(v) \text{ is \textbf{none} or \textsc{open}},\\
    \param{route\_direction\_penalty} & a(v) \text{ is \textsc{draining}},\\
    0 & \text{the traversal direction matches } a(v)\text{'s},\\
    \param{route\_direction\_penalty} & \text{otherwise,}
  \end{cases}\\
  \mathrm{cost}(u, v) &= 1 + \mathrm{penalty}(u, v),
\end{align}
with \param{route\_direction\_penalty} $= 6.0$ by default.

The rationale is stated in the module's own docstring: PIBT looks exactly one
step ahead, so directional information has to be baked into the \emph{route}
itself, or a robot will sit at an aisle entrance repeatedly proposing the same
badly-scoring move until the aisle flips. Note that this is a routing
preference, not a legality rule: an expensive edge is still an edge, in keeping
with \Cref{sec:contribution-soft}.
